\section{Methodology}
\label{sec:methodology}

\subsection{Measurement audit and withdrawn results}
\label{sec:audit}

% Before the controlled campaign, an internal audit examined the project's early
% wall-clock comparisons and found them unsuitable as performance evidence.
% The principal defects were: (i)~PDX and FAISS candidate timers excluded
% token-to-document aggregation, candidate selection, reranking, and final
% ranking; (ii)~PLAID was allowed up to 8{,}192 full scores per query while the
% IVF pipelines commonly reranked 50 documents, so the compared methods
% performed very different work; (iii)~the exact baseline was a per-document
% Python/NumPy loop rather than a compiled kernel; and (iv)~measurements from
% Windows and WSL environments appeared in the same figures. Additional issues
% included uncontrolled thread counts, no repeated or interleaved runs,
% parameter selection on the evaluation queries, favorable corpus prefixes, and
% IVF bucket counts that changed with corpus size in scaling plots.
% The historical headline ratios (approximately $30\times$, $39\times$, and
% $32\times$) were consequently withdrawn. Legacy artifacts are preserved under
% \texttt{results/legacy/} for mechanism and quality analysis only; no
% performance claim in this report relies on them.

Before the controlled campaign, an internal audit examined the project's early 
wall-clock comparisons and found them unsuitable as performance evidence. The 
principal defects were the following:
\begin{itemize}[noitemsep]
  \item PDX and FAISS candidate timers excluded token-to-document aggregation, 
  candidate selection, reranking, and final ranking.
  \item PLAID was allowed up to 8,192 full scores per query while the IVF 
  pipelines commonly reranked 50 documents, so the compared methods performed 
  very different work.
  \item The exact baseline was a per-document Python/NumPy loop rather than 
  a compiled kernel.
  \item Measurements from Windows and WSL environments appeared in the same figures.
\end{itemize}
Additional issues included uncontrolled thread counts, no repeated or interleaved 
runs, parameter selection on the evaluation queries, favourable corpus prefixes, 
and IVF bucket counts that changed with corpus size in scaling plots. The historical 
headline ratios (approximately $30\times$, $39\times$, and $32\times$) were 
consequently withdrawn. Legacy artifacts are preserved under \texttt{results/legacy/} 
for mechanism and quality analysis only; no performance claim in this report relies 
on them.

\subsection{Controlled benchmark protocol}
\label{sec:protocol}

All results in \cref{sec:results-ivf,sec:results-plaid,sec:results-bond,sec:results-pca} were
produced under the replacement protocol recorded in\\
\texttt{docs/benchmark\_protocol.md}, whose core requirements are:

\begin{itemize}[noitemsep]
  \item \textbf{Standardized setup.} All compared arms run in a single
    Python process on one machine (Ubuntu WSL2 on an AMD Ryzen~7~5800H),
    pinned to CPU affinity $\{0, 2, 4, 6\}$ (one hardware thread on each of
    four physical cores) with four OpenMP/BLAS threads.
  \item \textbf{Matched output and declared work contract.} Every arm must
    emit ranked document identifiers under the same exact MaxSim definition.
    IVF arms are compared at matched rerank budgets. PLAID's public API does
    not expose realized candidate counts, so its \texttt{n\_full\_scores}
    values are reported as configured upper-bound budgets rather than equated
    with observed IVF work.
  \item \textbf{Complete online timing boundary.} The online timer starts with
    normalized query token vectors resident in memory and stops when final
    ranked document identifiers are available. It covers index search,
    token-to-document mapping, aggregation, candidate selection, exact
    reranking, and final top-$k$ production. Query encoding and index
    construction are excluded and reported separately as offline costs.
    Reported ratios therefore describe this online boundary, not an
    end-to-end retrieval service.
  \item \textbf{Repetition schedule.} One untimed warm-up pass and five
    measured repetitions, with the arm order interleaved per repetition to
    reduce thermal and cache-order bias. Median and 95th-percentile (p95)
    latencies are reported and raw samples are preserved in the artifacts.
  \item \textbf{Provenance.} Every result JSON records the Git commit and
    dirty flag, input SHA-256 hashes, package versions, thread environment,
    CPU affinity, stage timings, and per-repetition samples. The release
    manifest accepts only artifacts from clean commit \texttt{0a11fda}.
\end{itemize}

\subsection{Qualitative metrics}
\label{sec:metrics}

Two families of qualitative metrics are kept strictly separate:

\begin{description}[noitemsep]
  \item[Exact-ranking recovery.] The primary architecture metric. For each
    query, the fraction of the \emph{exact} MaxSim top-$k$ documents that the
    approximate arm returns in its top-$k$ (reported at $k = 10$ as exact
    recall@10). Candidate-pool recovery measures the same quantity against the
    entire retrieved pool, isolating pool-coverage loss from selector loss.
  \item[Qrels metrics.] Standard relevance-judgment metrics (Recall@$k$,
    MRR@10) against the BEIR judgments. Because qrels are incomplete and the
    query samples are small, these metrics are insensitive to architecture
    differences (\cref{sec:results-ivf}) and are reported for completeness
    only.
\end{description}

\subsection{Datasets, splits, and configuration}

For each dataset the first 10 queries were used for pilot validation and
selector-policy freezing, and the remaining \textbf{40 held-out queries} form
the release comparison. On SciFact, the selector study found that a more
complex \texttt{union\_approx\_and\_count} policy changed exact recall@10 by
at most $\pm 0.02$ across budgets, so the simpler \texttt{approx\_score}
selector was frozen before the held-out runs. The candidate configuration is
fixed at $L = 100$ token hits per query token and $\texttt{nprobe} = 8$; only
the rerank budget $C \in \{50, 100, 200, 400, \text{full pool}\}$ varies
(SciFact; NFCorpus uses $\{50, 200, \text{full pool}\}$). ``Full pool'' means
that every unique document in the candidate pool is reranked.

PLAID parameters were selected only on the first 10 SciFact queries. The
validation sweep varied \texttt{n\_ivf\_probe} from 8 to 128 and
\texttt{n\_full\_scores} from 50 to 400. The fastest configuration reaching
validation exact recall@10 of at least 0.85 was $(8,100)$ (recall 0.89), and
the maximum-recall configuration was $(8,400)$ (recall 0.92). An expanded
validation-only sweep through 5{,}183 full scores did not exceed 0.92. The two
selected configurations were frozen before either held-out dataset was run.
After the unexpectedly large validation--test shift was observed, one
full-corpus \texttt{n\_full\_scores} arm per dataset was added as an explicitly
post-hoc sensitivity analysis; it is not treated as a selected operating point.

\subsection{Numerical precision contracts}
\label{sec:precision}

Two distinct precision contracts are used, and their timings are never mixed:
the \textbf{IVF comparison} reranks with the compiled float32 exact kernel,
matching the float32 ColBERT embeddings, whereas the \textbf{BOND comparison}
uses float64 products and accumulation in \emph{both} the exhaustive and the
BOND arm, because the exact-safe pruning proof requires controlled rounding
(\cref{sec:results-bond}). Consequently, float32 IVF latencies (e.g., the
\SI{2.58}{\second} SciFact exhaustive median) must not be compared with
float64 BOND-comparison latencies (e.g., the \SI{3.92}{\second} SciFact
exhaustive median); each table names its own precision-matched baseline, and
every BOND-to-exact ratio uses the exhaustive arm interleaved in the same run.
